\scanpage{82}
\deleted{
Cas d'un \deleted{\unsure{groupoïde algébrique}}\inserted{hom de groupoïdes}
\diagram{\begin{tikzcd}[ampersand replacement=\&]\widetilde{\mathcal T}\arrow[r,"p"]\&\mathcal T\end{tikzcd}}
qui est surjectif sur les $\pi_0$. On se donne une famille $(\widetilde u_i)_{i\in I}$ de flèches de $\widetilde{\mathcal T}$, \deleted{entre des objets} de projections $u_i$ dans $\mathcal T$ et des relations
\inserted{$\deleted{\unsure{i^{\alpha}}}\ i^{(\alpha)}=(i_0^{(\alpha)},\ldots,\unsure{i_{m_\alpha}^{(\alpha)}})\in I'^{m_\alpha}$}
\[(4)\qquad\unsure{u'_{i^{(\alpha)}_{m_\alpha}}}\,u'_{i^{(\alpha)}_{m_\alpha-1}}\cdots u'_{i^{(\alpha)}_1}u'_{i^{(\alpha)}_0}=1\quad\text{dans }\mathcal T\qquad(\alpha\in A)\]
\inserted{[pour le rôle des relations $R_\alpha$]}
(flèches composables dans $\mathcal T$, avec $\operatorname{or}(u'_{i_0})=\operatorname{ext}(u'_{i_{m-1}})$).
\margin{d'où une famille de flèches $\widetilde u'_{i'}$, $\deleted{\unsure{u_i}}=\widetilde u_i^{\pm1}$ paramétrée par les $i'\in I'\simeq I\times\{\pm1\}$, en posant $\widetilde u'_{(i,\varepsilon)}=\widetilde u_i^\varepsilon$ --- itou pour des $u'_{i'}$}
De plus, on se donne des \deleted{relations}\inserted{flèches} $(a_j)_{j\in J}$ de $\unsure{\mathcal T|\widetilde{\mathcal T}}$ qui sont des
\[\deleted{(5)\qquad\unsure{u}}\]
``flèches fibres'' (leur projection dans $\unsure{\widetilde{\mathcal T}|\mathcal T}$ des identités), \deleted{et des relations}\inserted{d'où des $a'_{j'}=a_j^{\pm1}\quad(j'\in J\times\{\pm1\})$}
\deleted{\unsure{S}} et des relations
\[(5)\qquad a'_{j^{(\beta)}_{n_\beta}}a'_{j^{(\beta)}_{n_\beta-1}}\cdots a'_{j^{(\beta)}_0}=1\]
\[\text{où }j^{(\beta)}=(j^{(\beta)}_0,\ldots,j^{(\beta)}_{n_\beta})\in J'^{\unsure{n_\beta}},\qquad\beta\in B\]
[pour le rôle des rel. $S_\beta$]
\margin{\uline{NB} Les indices \unsure{$i$ sont des éléments de} $\coprod_{n\in\mathbb N^*}I'^n$ \unsure{et $B$ un ens. d'éléments de} $\coprod_{n\in\mathbb N^*}J'^n$}
On suppose

1°) les relations (4) sont une famille \deleted{qui} de relations \unsure{définissantes} pour $\mathcal T$, en \inserted{flèches génératrices $u_i$ $(i\in I)$ et $\operatorname{Ob}\mathcal T$}
(ce qui implique que tous les sommets de $\mathcal T$ sont parmi les origines \inserted{et extr.} des flèches \deleted{\unsure{$u'_{i_j}$ intervenant dans (4)}} (à fortiori $\widetilde{\mathcal T}\text{---}\mathcal T$ est surjectif sur les objets). \deleted{et des}
}
\mspage{81}
